% Complete new receiving proof; original predecessor files remain preserved.
\section{The evaluated Gamma term in the original arithmetic and cohomology receiver}
\label{sec:WGR}

Retain the actual packet $k=4l+1$, $e=1+k(m-1)$,
$q=e(k+1)^2$, $c=k/2$, $0<\delta<1/2$, $\gamma>2$,
and the original polynomials
\[
 \chi(S)=\prod_{a,b=0}^k
 [S-c-(2a-k)\delta-i(2b-k)\gamma]^e,\qquad
 f(S)=\prod_{j=0}^{l-1}(S-c-2j-1/2).
 \tag{WGR1}
\]
Here $E=\mathbb C[S]/(\chi)$, $J_N:\mathbb C[S]_{\le N}\to E$
is the literal remainder map, and the source line remains $S=c+iy$.
The parameters, primary multiplicities, and polynomial maps have not
been changed by the evaluation of $W_k$.

Let $\mathfrak r_\chi$ be the original scalar norm ratio at $N=q$,
and let $\mathfrak R_{2q-1},\mathfrak R_{2q}$ be the two original
relation-Gram ratios from JSR12. Write the exact scalar errors
\[
 \begin{aligned}
 e_0&=\log\mathfrak r_\chi-\log(m_{2q+2l}/m_{2q}),\\
 e_r&=\log\mathfrak R_{q+r-1}
       -\log\frac{\det\mathsf M_{q+l,r}}{\det\mathsf M_{q,r}},
       \qquad r=q,q+1,\\
 (\mathsf M_{a,r})_{ij}&=\int_{\mathbb R}y^{2a+i+j}
            \frac{|\Gamma(1/4+iy/2)|^2}{2\pi}\,dy,
            \qquad 0\le i,j<r .
 \end{aligned}\tag{WGR2}
\]
Thus the exponent index $m_j$ is the literal power $j$, and both
matrices use the full original mass $\sqrt{2\pi}$.
The exact coefficient map $P(S)\mapsto P(c+iy)$ has entries
$J_{hj}=\binom jh c^{j-h}i^h$ and determinant of absolute value one.
In particular the high reference Gram is its congruence
$J^*\mathsf M_{a,r}J$, not a different-coordinate determinant with
an unrecorded power of $q$. The bulk multiplication
$\chi(c+iy)/(iy)^q$ and the full positive inner and far forms are
the exact source maps and corrections proved in CTR and LET.

On the explicit domain $k\ge2048\sqrt{\delta^2+\gamma^2}$,
LET18 and CTR48 prove
\[
 e_0\in[-a_0,b_0],\qquad e_q+e_{q+1}\in[-A,B],
 \quad A=2lE_0+2(L_q+L_{q+1}),\quad
 B=A+(2q+1)\Delta_f,
 \tag{WGR3}
\]
where $a_0=E_0+2L_1$, $b_0=a_0+\Delta_f$, and every constant
$E_0,\Delta_f,L_r$ is the finite original integral constant printed
in LET7--16 and BHR2. In particular $A,B=O(k)$ for fixed $m=1$
and $O(1)$ for fixed $m\ge2$; $a_0,b_0=O(1)$ and $O(k^{-1})$
in those respective branches. These statements were proved on the
full source, including the low rank one, and are used within their
stated domain here.

The exact identity LET19, retaining the low minus sign, is
\[
 \mathfrak T_k=W_k-e_0+e_q+e_{q+1},\qquad
 \mathcal H_k:=R_k^0-R_k^\sigma=-\mathfrak T_k.
 \tag{WGR4}
\]
RWB20--27 now supplies an evaluated pair of finite numbers
\[
 \begin{aligned}
 w^-_{q,l}=\max\bigg\{\frac{lq}{64},\;&
 4lq\log\frac{27}{8\pi}-4l^2-6l\log(3/2)-3l\\
 &-\frac{9l}{4q-1}-\frac{2l(l+1/4)}q\bigg\},
 \qquad w^+_{q,l}=6lq,\qquad
 w^-_{q,l}\le W_k\le w^+_{q,l}.
 \end{aligned}\tag{WGR5}
\]
This assertion holds at every original $l,m\ge1$; the extra threshold
in WGR3 is required only when returning to the original $\chi$ source.
Substituting WGR5 and then the three signed errors into WGR4 proves
\[
 \boxed{-w^+_{q,l}-B-a_0\le\mathcal H_k
                 \le-w^-_{q,l}+A+b_0.}
 \tag{WGR6}
\]
Indeed $\mathcal H_k=-W_k+e_0-e_q-e_{q+1}$.
This proves both signs separately; no positive high error has been
assigned the sign of the negative low endpoint.

\paragraph{The complete retained arithmetic and mixed-state map.}
Use exactly the previous signed arithmetic difference
$\delta_k^\sigma$ and the full baseline $\mathcal B_k^0$ from JSR16--21.
For either original source, the quotient, kernel and boundary forms
are
\[
 G_N=(J_NH_N^{-1}J_N^*)^{-1},\qquad
 H_N^K=I_K^*G_NI_K,\qquad
 Q_N^B=(\Lambda G_N^{-1}\Lambda^*)^{-1}.
 \tag{WGR7}
\]
Here $\Lambda:E\to B_{\mathrm{obs}}=\operatorname{im}\Lambda$ is the
original boundary observation, $K=\ker\Lambda$, and $I_K:K\hookrightarrow E$
is the original inclusion. The chosen original lift
$S_*:B_{\mathrm{obs}}\to E$ satisfies $\Lambda S_*=I_{B_{\mathrm{obs}}}$.
Thus $[I_K,S_*]:K\oplus B_{\mathrm{obs}}\to E$ is the exact
coefficient isomorphism. The exact sequence is represented in this common
coefficient frame $[I_K,S_*]$. Orthogonal projection subtracts a kernel
column from each lift in $S_*$ and is a triangular change with
determinant one, proving
\[
 \det H_N^K\det Q_N^B
       =|\det[I_K,S_*]|^2\det G_N.
 \tag{WGR8}
\]
The same frame is used for both sources, so its determinant cancels
in their ratio. The source comparison gives the original nonnegative
deficits $b_N^K+b_N^B=b_N\le\Xi_{k,N}$ and the four signed sums
\[
 \begin{aligned}
 X_k&=-b_{q-1}-b_q+b_{2q-1}+b_{2q},\\
 a_k^-&=\Xi_{k,q-1}+\Xi_{k,q},\quad
 a_k^+=\Xi_{k,2q-1}+\Xi_{k,2q},\qquad
 -a_k^-\le X_k\le a_k^+ .
 \end{aligned}\tag{WGR9}
\]
The same formula applies to $\Delta F_K,\Delta F_B$ by using the
corresponding superscripts, and to $\delta_k^\sigma$ by using the
full quotient. The last inequality follows by dropping the two
positive terms for the lower bound and the two negative terms for
the upper bound. Thus the original mixed kernel and boundary
directions remain in the exact map, rather than being lost in a
scalar positivity assertion. The previous finite source estimates
$a_k^-\le D_k^-$ and $a_k^+\le D_k^+$, with all original constants
in JSR17--18, continue to apply separately.

The original arithmetic equality and the evaluated finite bounds are
therefore
\[
 \begin{aligned}
 \mathcal B_k^{\rm ar}&=\mathcal B_k^0+\mathcal H_k+\delta_k^\sigma,
       &\Delta_k^\Gamma&=\mathcal H_k+\delta_k^\sigma,\\
 \mathcal B_k^0-w^+_{q,l}-B-a_0-a_k^-
 &\le\mathcal B_k^{\rm ar}
       \le\mathcal B_k^0-w^-_{q,l}+A+b_0+a_k^+ .
 \end{aligned}\tag{WGR10}
\]
Add WGR9 to WGR6 to prove WGR10. Replacing $a_k^\pm$ by their
respective outer bounds $D_k^\pm$ gives the corresponding fully
estimated arithmetic interval. In particular the original nonnegative
HC residual has the finite enclosure
\[
 \begin{aligned}
 \max\{0,\mathcal B_k^0-w^+_{q,l}-B-a_0-D_k^-
                     -4q\log(D_hk)\}
 &\le\mathcal R_k:=\mathcal B_k^{\rm ar}-4q\log(D_hk)\\
 &\le\mathcal B_k^0-w^-_{q,l}+A+b_0+D_k^+
                     -4q\log(D_hk).
 \end{aligned}\tag{WGR11}
\]
Its constant $D_h$ and its threshold are the previous original ones.
For every finite exact $W_k$, WGR6 contains the sharper WRC5 interval.
Consequently it provides an evaluated bound with the original source
constants on the growing family; it does not claim to narrow an interval already using
the exact scalar $W_k$ as input.

\paragraph{A definite growing arithmetic contribution.}
Fix an original packet $(\delta,\gamma,m)$ and let $k=4l+1\to\infty$.
The already proved JSR18--21 finite source estimates give
\[
 \begin{cases}
 \delta_k^\sigma/(kq)=O_h((\log k/k)^{5/7}),&m=1,\\
 \displaystyle\liminf\frac{\delta_k^\sigma}{q\log k}\ge-(64m+245),
 \quad\displaystyle\limsup\frac{\delta_k^\sigma}{q\log k}
                                      \le128m+490,&m\ge2.
 \end{cases}\tag{WGR12}
\]
Thus $\delta_k^\sigma/(kq)\to0$ in both branches. In the second
branch multiply the finite upper and lower constants by $\log k/k\to0$;
in the first branch the displayed rate tends to zero. Also WGR3 implies
$(e_0-e_q-e_{q+1})/(kq)\to0$, retaining its original asymmetric bounds.
Since $l/k=(k-1)/(4k)\to1/4$, RWB23 and RWB27 yield
\[
 \boxed{-\frac32\le\liminf\frac{\Delta_k^\Gamma}{kq}
 \le\limsup\frac{\Delta_k^\Gamma}{kq}
 \le-\log\frac{27}{8\pi}<0.}
 \tag{WGR13}
\]
This uses $\Delta_k^\Gamma=-W_k+e_0-e_q-e_{q+1}+\delta_k^\sigma$
and the proved finite bounds, so it does not infer a limit from the
exploratory numerical recurrence. In particular the Gamma/arithmetic
difference is strictly negative for all sufficiently large original
degrees. It remains a term in the full arithmetic equality WGR10.

The preceding full-source equilibrium calculation GEL18--19 evaluates
this term more precisely, proving
\[
 \boxed{\frac{\Delta_k^\Gamma}{kq}\longrightarrow
       -\frac{\mathcal C_\Gamma}{4},\qquad
 4\log\frac{27}{8\pi}\le\mathcal C_\Gamma\le6.}
 \tag{WGR14}
\]
Indeed GEL18 proves $W_k/(lq)\to\mathcal C_\Gamma$, while the
two other summands divided by $kq$ vanish by WGR12 and WGR3.
The finite RWB bounds give the stated interval for the exact equilibrium
integral $\mathcal C_\Gamma$. Combining this result with WGR10--11
gives the evaluated relation to the original baseline and residual:
\[
 \frac{\mathcal B_k^0-\mathcal B_k^{\rm ar}}{kq}
       \longrightarrow\frac{\mathcal C_\Gamma}{4},\qquad
 \frac{\mathcal R_k-\mathcal B_k^0}{kq}
       \longrightarrow-\frac{\mathcal C_\Gamma}{4},\qquad
 \liminf\frac{\mathcal B_k^0}{kq}\ge\frac{\mathcal C_\Gamma}{4}.
 \tag{WGR15}
\]
The second assertion also uses $4\log(D_hk)/k\to0$ at fixed $D_h>0$;
the third follows from the independent original inequality
$\mathcal R_k\ge0$. These are evaluated differences and a necessary
baseline bound with the original arithmetic object retained. They do
not assert that the baseline or the nonnegative residual has vanished.
The finite $q$-scale terms in WGP13 and the finite allowances in
WGR10--12 remain available for the next calculation on those actual
objects. In particular the $o(lq)$ remainder in GEL18 has not been
replaced by an unproved $o(q)$ remainder.
